% ============================================
% Supplementary Note 1: Robustness of Coalescence Outcome Classification and Null-Model Controls
% ============================================

\subsection*{Supplementary Note 1: Robustness of coalescence outcome classification and null-model controls}
\label{sec:suppl:note1}

The classification framework, similarity metric, and outcome thresholds used throughout this work are defined in Supplementary Methods. Here we test whether the predominance of Dominance is robust to these analytical choices and whether it could instead arise as an artifact of classification boundaries, compositional metric, abundance skew, or a few dominant species.

\paragraph{Boundary sensitivity of the classification thresholds.}
To make the boundary dependence explicit, we report the continuous similarity measures underlying the categorical classifications across media in Supplementary Figs.~10--12. A boundary-sensitivity analysis (Supplementary Fig.~37) further shows that the nutrient-dependent increase in Dominance remains evident when either the PDI boundary or retention boundary is varied around the manuscript baseline.

\paragraph{Abundance-skew null models.}
A potential concern is that Dominance could arise from skewed abundance distributions rather than ecological interactions. To address this, we compared experimental parental-asymmetry values against two null models.

In Null Model 1 (abundance-weighted random selection), species survival probability is proportional to abundance in the combined parental-community pool, testing whether high-abundance species simply survive regardless of parental-community affiliation. In Null Model 2 (shuffled abundance), abundances are randomly permuted among species within each parental community before neutral mixing, testing whether the overall skewness structure alone drives Dominance.

We generated 500 null offspring communities per model and calculated parental asymmetry, $|2\mathrm{PDI}-1|$ (Supplementary Fig.~44). Experimental Base-medium data showed mean parental asymmetry of 0.812 ($n = 83$). Both null models produced significantly lower values: abundance-weighted null mean = 0.503 ($U = 32188$, $p = 7.93 \times 10^{-16}$), shuffled null mean = 0.552 ($U = 32050$, $p = 1.85 \times 10^{-15}$; Mann-Whitney U tests). These results indicate that abundance skewness alone cannot explain the observed patterns and support the conclusion that ecological processes beyond simple abundance-weighted survival help determine which parental community dominates.

\paragraph{Simple additive model.}
We also tested a simple additive null model, $n_{C,\mathrm{null}} = n_A+n_B$, constructed separately for each experimental coalescence event from the raw ASV count vectors of that event's two parental communities (Extended Data Fig.~3). Under this model, an apparent Dominance classification can arise when the two parental-community raw ASV count vectors have substantially different Euclidean norms: for non-overlapping parental-community raw ASV count vectors, $\mathrm{Sim}(A,A+B)$ and $\mathrm{Sim}(B,A+B)$ are proportional to $\|n_A\|_2$ and $\|n_B\|_2$, respectively. This norm-imbalance effect is possible in principle, especially when richness is low and vector norms are sensitive to abundance unevenness among a small number of taxa. In the experimental data, however, the simple additive null model was usually classified as Mixture, including 66/83 Base medium events (79.5\%), and the observed Base medium outcomes were shifted further toward parental asymmetry than their corresponding null-model outcomes (one-sided paired Wilcoxon test, $p = 5.9 \times 10^{-14}$). Thus, the simple additive null model does not explain the observed Dominance tendency. Related geometric effects can arise more generally when parental-community composition vectors differ in norm or unevenness; these are evaluated both with the simple additive null model above and with the simulation parental-community composition-vector norm analysis in Supplementary Fig.~38. Together, these analyses show that such geometric effects can produce apparent Dominance in principle, but do not account for most experimental coalescence outcomes or the simulated Dominance transition.

\paragraph{Sensitivity of outcome classification to the compositional metric and its biological interpretation.}
A further concern is whether the predominance of Dominance is an artifact of the particular compositional metric used to define the outcome categories rather than a robust feature of the data. To test this, we compared coalescence outcome classifications using five approaches: the primary cosine-similarity/vector-decomposition framework, Euclidean distance, Bray-Curtis dissimilarity, Jensen-Shannon divergence, and Jaccard index. The alternative relative-abundance approaches, Euclidean distance and Bray-Curtis dissimilarity, produced outcome distributions broadly consistent with the primary framework, with Dominance as the most frequent outcome. However, Jensen-Shannon divergence and Jaccard index did not recover this pattern (Extended Data Fig.~2), reflecting the sensitivity of boundary assignments to the choice of compositional metric. We interpret this divergence as biologically informative rather than solely technical: the primary framework, Euclidean distance, and Bray-Curtis retain quantitative abundance information, whereas Jaccard collapses the data to presence/absence and therefore measures species-identity retention rather than abundance-weighted parental dominance. The Dominance pattern is therefore robust across abundance-weighted parental-similarity metrics. Its absence under Jaccard indicates that Dominance is not primarily a presence/absence retention signal, while the Jensen-Shannon result indicates that full-distribution divergence emphasizes a different aspect of community composition than directional parental similarity.

\paragraph{Effect of dominant-species abundance on PDI and the Dominance trend.}
Because PDI is abundance-weighted, a single highly abundant taxon shared between the coalesced community and one parental community can dominate the composition vectors and inflate the directional similarity toward that parent. This raises the concern that the nutrient-dependent increase in Dominance (Fig.~4D) and the Fig.~5C predictability analysis could partly reflect a mechanical contribution from a few dominant species rather than a community-level reorganization toward one parental community---that is, the association could be partly circular, because the same dominant taxa that define the coalesced community also determine its PDI. To test this, we recalculated PDI after removing either the coalesced-community dominant species or the two parental-community dominant species from all three composition vectors, so that the metric reflects only the subdominant community structure. The Nutr$+$ association weakened but remained significant, whereas Base was not significant after removal (Supplementary Fig.~16). Thus, the Nutr$+$ Dominance signal is partly, but not entirely, attributable to the most abundant species; a community-wide contribution beyond the dominant taxa remains.
